\section{Introdução}

A transformação digital tem exigido das organizações não
apenas a adoção de novas ferramentas, mas a revisão de seus
processos internos, especialmente aqueles que ainda dependem
de comunicação informal entre setores. Quando problemas
recorrentes são reportados por aplicativos de mensagens ou
e-mails não estruturados, sem protocolo e sem registro
centralizado, a organização perde rastreabilidade, dificulta
a gestão e impossibilita a geração de indicadores de
desempenho \cite{Ramos2024}. A literatura recente sobre
transformação digital é enfática ao afirmar que a
digitalização de um fluxo manual, quando realizada sem a
revisão estrutural do processo que lhe deu origem, tende
apenas a reproduzir em ambiente eletrônico as mesmas
ineficiências já existentes \cite{SantosCaldeiraPacheco2026}.
Este trabalho parte dessa premissa: antes de desenvolver uma
aplicação, é necessário compreender e redesenhar o processo
que ela pretende sustentar.

Assim, a lógica de construção desta pesquisa segue uma
sequência específica --- do problema organizacional à
automação --- que também organiza a estrutura dos capítulos
deste trabalho: identificação do problema organizacional,
análise e entendimento do processo existente, identificação
dos gargalos e necessidades, levantamento de requisitos,
redesenho do processo e, somente então, desenvolvimento da
solução tecnológica e automação. Essa sequência é sustentada
pela integração entre a Engenharia de Processos, o
\textit{Business Process Management} (BPM), a Engenharia de
Software, a Engenharia de Requisitos e a abordagem
\textit{low-code}, discutidas na Base Teórica deste trabalho.

Em organizações que ainda não dispõem de um canal único de
abertura de chamados internos, problemas que exigem análise e
resolução por diferentes áreas são comunicados de forma não
padronizada. Não há protocolo gerado automaticamente, não há
controle de prazo e não há visibilidade gerencial sobre o
volume, a prioridade ou o tempo médio de resolução dessas
demandas. A ausência de um canal formalizado de entrada expõe
a organização ao retrabalho, à perda de informações e à
impossibilidade de identificar causas raiz dos problemas
recorrentes \cite{Santos2023}.

Diante desse cenário, este trabalho realiza a análise do
processo atual (AS-IS), o levantamento de requisitos, a
modelagem do processo futuro (TO-BE), o desenvolvimento e a
implantação de um fluxo de gestão de chamados internos entre
áreas, tendo como ponto de entrada um site desenvolvido no
\textbf{SharePoint}, responsável por centralizar a abertura
das solicitações. Os chamados registrados nesse site
constituem a entrada de um processo automatizado desenvolvido
com base no ecossistema Microsoft~365, contemplando o registro
estruturado da solicitação, a padronização do fluxo entre as
áreas envolvidas e a geração de indicadores gerenciais de
acompanhamento \cite{Lopes2024}.

O objetivo geral deste trabalho é analisar o processo atual
de comunicação de chamados internos entre áreas, levantar os
requisitos, redesenhar e modelar o processo, e desenvolver e
implantar um sistema de gestão de chamados internos, tendo
como ponto de entrada um site no SharePoint e utilizando
ferramentas do ecossistema Microsoft 365, de modo a
padronizar a abertura das solicitações, eliminar a
comunicação informal e viabilizar a geração de indicadores
gerenciais de desempenho do processo.

% ─── JUSTIFICATIVA ────────────────────────────────────────
\subsection{Justificativa do Tema}

A escolha do tema fundamenta-se em quatro dimensões
complementares: relevância prática imediata, aderência à
disciplina de Engenharia de Processos e de Gestão de
Processos de Negócio, aderência às competências da
Engenharia de Software e da Engenharia de Requisitos, e
contribuição ao conhecimento acadêmico aplicado.

Do ponto de vista prático, a organização estudada não dispõe
de qualquer mecanismo formalizado para abertura,
acompanhamento e encerramento de chamados internos entre
áreas. O autor possui conhecimento direto das dificuldades
operacionais do processo e acesso às informações necessárias
ao seu diagnóstico e ao levantamento de requisitos
\cite{Gil2002}. A ausência de um canal formalizado de abertura
de chamados gera custos ocultos decorrentes de retrabalho,
perda de informações em canais informais e impossibilidade de
gestão baseada em dados.

Do ponto de vista da Engenharia de Processos e do BPM, o tema
exige que a solução tecnológica seja precedida pela
compreensão do processo existente e pelo seu redesenho
estruturado, superando a lógica de que a automação, por si
só, resolve problemas de gestão \cite{SantosCaldeiraPacheco2026}.
Do ponto de vista da Engenharia de Software e da Engenharia
de Requisitos, o tema mobiliza competências nucleares do
curso: especificação de requisitos a partir da modelagem do
processo de negócio \cite{CarvalhoJunior2022}, arquitetura de
solução orientada a serviços e avaliação de qualidade com
usuários. A construção de um sistema de chamados com controle
de status, prioridade, SLA, causa raiz e rastreabilidade,
integrado a um painel gerencial, constitui um projeto de
engenharia completo, com requisitos funcionais e não
funcionais bem definidos, mesmo em ambiente \textit{low-code}
\cite{Alves2022}.

Do ponto de vista acadêmico, a pesquisa contribui com
evidência empírica sobre a aplicação integrada da Engenharia
de Processos, do BPM, da Engenharia de Requisitos e do
\textit{low-code} na modelagem de um processo de gestão de
chamados internos, a partir de um site no SharePoint como
ponto único de entrada, integrado ao ecossistema Microsoft
365 --- recorte pouco explorado de forma conjunta na
literatura nacional recente, constituindo contribuição
original ao campo \cite{Santos2023}.

% ─── PROBLEMA ─────────────────────────────────────────────
\subsection{Problema da Pesquisa}

Os chamados internos entre áreas da organização são
comunicados por aplicativos de mensagens e e-mails sem
padronização, sem geração de protocolo, sem controle de prazo
e sem registro centralizado. A área responsável pela análise
recebe a demanda sem garantia de completude das informações.
A gestão não possui indicadores sobre volume, prioridade,
tempo de resolução ou causa raiz dos problemas. Chamados são
esquecidos ou duplicados sem que haja mecanismo de alerta ou
escalonamento.

Diante desse contexto, formula-se a seguinte questão de
pesquisa:

\textit{Como estruturar, a partir da análise do processo
atual e de um site no SharePoint como ponto único de entrada,
o levantamento de requisitos, o redesenho, a modelagem, o
desenvolvimento e a implantação de um processo automatizado
de gestão de chamados internos entre áreas, capaz de
padronizar a abertura, garantir rastreabilidade, controlar
SLA por prioridade, registrar causa raiz e gerar indicadores
gerenciais de volume, tempo médio de resolução e percentual
de chamados concluídos dentro do prazo?}

% ─── HIPÓTESE ─────────────────────────────────────────────
\subsection{Hipótese}

A análise do processo atual, o levantamento de requisitos, o
redesenho segundo os princípios do BPM, o desenvolvimento e a
implantação de um processo de gestão de chamados internos
entre áreas --- tendo como ponto de entrada um site no
SharePoint e como componentes de automação a Microsoft List,
o Power Automate e o Power BI --- permitem obter uma solução
capaz de eliminar a comunicação informal entre setores,
garantir rastreabilidade integral dos chamados, controlar SLA
por prioridade, registrar causa raiz e ação corretiva, e
gerar indicadores gerenciais de volume, tempo médio de
resolução e percentual de atendimentos dentro do prazo,
contribuindo para a melhoria mensurável da eficiência
operacional das áreas envolvidas.

% ─── ESTADO DA ARTE ───────────────────────────────────────
\subsection{Estado da Arte}

O levantamento realizado para este trabalho abrangeu
produções em língua portuguesa, publicadas nos últimos cinco
anos e indexadas no Google Scholar, SciELO e Portal de
Periódicos CAPES, com os descritores: \textit{engenharia de
processos}, \textit{BPM}, \textit{engenharia de requisitos},
\textit{elicitação de requisitos}, \textit{low-code},
\textit{Microsoft List}, \textit{Power Automate},
\textit{automação de processos}, \textit{SLA} e
\textit{indicadores de desempenho}.

No campo da Engenharia de Processos e da transformação
digital, \cite{SantosCaldeiraPacheco2026} demonstram que a
transformação digital efetiva depende da revisão e do
redesenho estruturado dos processos organizacionais, e não da
simples informatização de fluxos preexistentes, resultado que
orienta diretamente a decisão deste trabalho de tratar o
redesenho do processo como etapa anterior ao desenvolvimento.
No campo da gestão de processos, \cite{Ramos2024} demonstram
que a disponibilização de dados gerenciais estruturados é o
resultado de maior valor percebido após iniciativas de BPM em
instituições de médio e grande porte. Em linha com esse
resultado, \cite{Lopes2024} evidenciam que a digitalização de
processos com BPM e ferramentas do ecossistema Microsoft
resultou em repositório único de dados com rastreabilidade
completa, eliminando a comunicação informal entre setores.

No campo da Engenharia de Requisitos, \cite{CarvalhoJunior2022}
demonstram que a derivação de requisitos de software a partir
da modelagem prévia do processo de negócio produz requisitos
mais alinhados às metas organizacionais do que técnicas
tradicionais de elicitação isoladas, resultado que reforça a
opção metodológica deste trabalho de conduzir o levantamento
de requisitos junto a perfis que representam diferentes níveis
de atuação no processo.

No campo das plataformas \textit{low-code},
\cite{Alves2022} identificam redução de até 60\% no tempo de
desenvolvimento de soluções corporativas com o paradigma
\textit{low-code} em comparação ao desenvolvimento
tradicional, destacando o ecossistema Microsoft 365 como
plataforma de maior penetração no mercado nacional.
\cite{Santos2023} concluem que a combinação de BPM com
\textit{low-code} é especialmente eficaz na construção de
repositórios e painéis de acompanhamento, pois permite que as
regras de negócio sejam configuradas pelas próprias áreas
proprietárias dos processos. \cite{Costa2022} demonstram, em
estudo de caso aplicado, que a adoção de ferramentas
\textit{low-code} resultou em rastreabilidade completa e
redução de 45\% no tempo de processamento de solicitações
internas.

No campo da automação de processos, \cite{ReisFaria2024}
evidenciam, em revisão sistemática sobre automação robótica
de processos, que os ganhos de padronização e redução de
retrabalho são mais consistentes quando a automação é
aplicada sobre um processo já revisado. No campo da gestão de
SLA, \cite{Malheiros2024} concluem que o SLA é fundamental
para estabelecer expectativas claras entre áreas, promovendo
qualidade, transparência e rastreabilidade na resolução de
chamados, e destacam que a automação do controle de prazos é
condição para que o SLA deixe de ser apenas um documento
formal.

O levantamento evidencia uma lacuna na literatura nacional
quanto à aplicação integrada da Engenharia de Processos, do
BPM e da Engenharia de Requisitos na construção de processos
de gestão de chamados internos entre áreas a partir do
ecossistema Microsoft 365, com um site no SharePoint como
ponto único de entrada. O presente trabalho busca preencher
essa lacuna com a análise do processo atual, o levantamento
de requisitos, o redesenho, o desenvolvimento e a implantação
do processo, documentados de forma estruturada.

% ─── OBJETIVOS ────────────────────────────────────────────
\subsection{Objetivos}

\subsubsection{Objetivo Geral}

Analisar o processo atual de gestão de chamados internos
entre áreas, levantar os requisitos, redesenhar e modelar o
processo, e desenvolver e implantar um sistema de gestão de
chamados internos, tendo como ponto de entrada um site no
SharePoint e utilizando plataformas do ecossistema Microsoft
365, de modo a padronizar a abertura, garantir
rastreabilidade, controlar SLA por prioridade, registrar
causa raiz e gerar indicadores gerenciais de desempenho do
processo.

\subsubsection{Objetivos Específicos}

\begin{itemize}
    \item Analisar o processo atual (AS-IS) de comunicação de
    chamados internos entre áreas, identificando participantes,
    atividades manuais e pontos de falha;

    \item Levantar e especificar os requisitos funcionais e
    não funcionais do processo por meio de um formulário
    estruturado, aplicado aos profissionais responsáveis pelo
    processo em diferentes níveis de atuação (Analista,
    Supervisora e Gerente), com perguntas direcionadas
    conforme o nível de responsabilidade de cada participante;

    \item Redesenhar o processo (TO-BE) a partir dos
    requisitos levantados, definindo os campos obrigatórios de
    abertura no site do SharePoint, os status possíveis
    (Aberto, Em análise, Concluído), os critérios de
    prioridade (Alta, Média, Baixa) e o tempo de SLA associado
    a cada nível de prioridade;

    \item Especificar o repositório centralizado a ser
    integrado ao site do SharePoint, com os campos: Protocolo
    N\textsuperscript{o}, Tipo de chamado (Massivo ou Interno),
    Assunto, Problema, Quantidade de protocolos relacionados,
    Área Responsável, Prioridade, Evidências, Solicitante,
    Status, Causa raiz, Ação corretiva, SLA (contagem
    automática) e Data da solução;

    \item Desenvolver o fluxo automatizado responsável por
    registrar cada chamado a partir do site do SharePoint,
    gerar protocolo único, enviar notificação à área
    responsável com link de acesso direto ao chamado e
    realizar a contagem automática do SLA;

    \item Desenvolver e implantar o painel gerencial com os
    indicadores: chamados abertos, em análise e concluídos por
    período; tempo médio de resolução; percentual de chamados
    dentro do SLA; distribuição por área responsável e por
    prioridade; e base histórica com rastreabilidade completa
    por protocolo;

    \item Validar o novo modelo do processo junto aos
    participantes envolvidos, a partir do fluxo implantado e
    do painel gerencial em uso.
\end{itemize}
